% !TeX root = ./main.tex

% --------------------------------------------------
% 資訊設定（Information Configs）
% --------------------------------------------------
% TODO(author): 以下個人資料皆為佔位符，請於提交前填入正確內容：
%   author / author* / ID / advisor / advisor* / college / institute /
%   oral-date / DOI，並確認中文題目翻譯。
% --------------------------------------------------

\ntusetup{
  university*   = {National Taiwan University},
  university    = {國立臺灣大學},
  college       = {電機資訊學院},                                
  college*      = {College of Electrical Engineering and Computer Science}, 
  institute     = {資訊工程學研究所},                          
  institute*    = {Department of Computer Science and Information Engineering}, 
  title         = {結合設計規劃、IP 重用與閉迴路修復的代理式 Verilog 程式碼生成框架}, 
  title*        = {An Agentic Framework for RTL Code Generation with Design Planning, IP Reuse, and Closed-Loop Repair},
  author        = {洪聿鍇},                                
  author*       = {Yu-Kai Hung},                             
  ID            = {R13922149},                          
  advisor       = {洪士灝},                           
  advisor*      = {Shih-Hao Hung},                         
  date          = {2026-07-17},                                  % TODO: 論文完成日期
  oral-date     = {2026-07-30},                                  
  DOI           = {10.6342/NTU202601985},                        % TODO
  keywords      = {大型語言模型, 暫存器傳輸級設計, 代理式規劃, IP 重用, 閉迴路修復, 模型路由, 硬體描述語言},
  keywords*     = {Large Language Models, RTL Generation, Agentic Planning, IP Reuse, Closed-Loop Repair, LLM Routing, Verilog},
}

% --------------------------------------------------
% 加載套件（Include Packages）
% --------------------------------------------------

\usepackage[numbers,sort&compress]{natbib}      % 參考文獻（數字式引用，配合 abbrv）
\usepackage{amsmath, amsthm, amssymb}   % 數學環境
\usepackage{ulem, CJKulem}              % 下劃線、雙下劃線與波浪紋效果
\usepackage{booktabs}                   % 改善表格設置
\usepackage{multirow}                   % 合併儲存格
\usepackage{diagbox}                    % 插入表格反斜線
\usepackage{array}                      % 調整表格高度
\usepackage{longtable}                  % 支援跨頁長表格
\usepackage{paralist}                   % 列表環境
\usepackage{listings}                   % 程式碼（Verilog / JSON 片段）
\usepackage{xurl}                       % 參考文獻長網址允許任意斷行
\usepackage{tikz}                       % 路由器流程圖
\usetikzlibrary{arrows.meta, positioning, shapes.geometric}

% --------------------------------------------------
% 套件設定（Packages Settings）
% --------------------------------------------------

% 放寬段落斷行（Times 字型下避免 Overfull hbox）
\emergencystretch=3em

% Verilog / generic code listings
\lstdefinelanguage{Verilog}{
  morekeywords={module, endmodule, input, output, inout, wire, reg, logic,
    assign, always, always_comb, always_ff, posedge, negedge, begin, end,
    if, else, case, casez, endcase, parameter, localparam, function,
    endfunction, generate, endgenerate, integer, initial},
  sensitive=true,
  morecomment=[l]{//},
  morecomment=[s]{/*}{*/},
  morestring=[b]",
}
\lstset{
  basicstyle=\ttfamily\small,
  keywordstyle=\bfseries,
  commentstyle=\itshape\color{gray},
  breaklines=true,
  columns=fullflexible,
  frame=single,
  framerule=0.3pt,
  numbers=none,
  showstringspaces=false,
  captionpos=b,
}
